\section{Extended social chronologies}
\label{app:social-detail}

\paragraph{\game{11}: a transient echo, a check and a rejection.}
At \decision{11}{255} Sol (seat~4) offers Virginia for Grok's (seat~2) mortgaged Pennsylvania plus \$150, publicly asserting ``Virginia completes your pinks''. Grok holds States but not St.~Charles, which Sol owns, so the claim is false. Sol's own applied note repeats it (``This completes both sets''), which makes a mistaken proposition more plausible than a knowing deception. Grok's returned reasoning at \decision{11}{256} first echoes the premise (``accepting would complete my pink monopoly'') and, a few hundred characters later, corrects it (``would not complete''). Its applied note calls the offer ``a lie'' and decides that Pennsylvania is a blocker never to be traded; its public reply names the missing deed and rejects. At \decision{11}{257} Sol concedes (``You're right---I missed St.~Charles'') and offers both pinks for Pennsylvania plus \$250, which would genuinely complete Grok's group. Grok rejects that too at \decision{11}{258}, stating that the deal would complete Sol's greens. Three observations follow. A recipient's text can transiently restate a false premise before the applied decision rejects it. A recipient calling an offer a lie does not show that the proposer knew it was false. And in this case own-group completion was not enough to buy a blocker when the same trade armed a rival.

\paragraph{\game{4}: public concession while the private record already held the true figure.}
In game~4, after Sol (seat~3) completes the pink group, Gemini (seat~4) ends its turn at \decision{4}{138} with a public invitation: ``Now that p3 has the Magenta set, p1 and p2---if either of you wants to work out a multi-way or direct deal so p3 isn't the only one with a monopoly, let me know.'' Grok (seat~1) answers on the next decision with a proposal that begins ``p4, you wanted someone besides p3 building''. To make a yellow set tradable, Grok first buys mortgaged Ventnor Avenue from Sol, the leader it is organizing against, bargaining from \$100 to \$200 between \decision{4}{142} and \decision{4}{146}; its notes weigh how many ``p3 pink houses'' each price would fund. The orange--yellow negotiation with Gemini then runs from \decision{4}{147} to \decision{4}{152}. At \decision{4}{148} Gemini's applied note records the true \$150 yellow house cost and a contingency: if Grok insists on taking orange, Gemini will ``demand cash/equity compensation''. At \decision{4}{149} Grok refuses, publicly claiming that yellow would leave it ``a yellow set I can't house---\$200 a house''. At \decision{4}{150} Gemini concedes publicly (``Fair point on house costs'') and switches to taking yellow with a \$143 cash request; its applied note justifies the switch by board position, and its returned reasoning retains the \$150 figure. Grok counters with \$50 at \decision{4}{151}, and Gemini accepts at \decision{4}{152}. The trade completes a group for both sides. Public assent to a false premise is therefore not sufficient evidence that the premise was adopted: the recipient's direction change matches a contingency it recorded before the false figure was spoken. Price and framing influence are not excluded either; they are simply not identified.

\paragraph{Bluffs that were bargaining posture.}
Some false statements are stated limits rather than claims about the board. At \decision{7}{42} Grok publicly calls \$600 for North Carolina Avenue ``final''; its applied note reads ``If they return at \$625, accept \ldots\ Do not go above \$625.'' Sol accepts \$600 at \decision{7}{43}, citing its own reason: to ``reinvest in cheaper, stronger sets rather than hold out for another \$25''. This is the closest case in the reviewed material of a recorded bluff followed by settlement at the bluffed figure, and the recipient's stated reason does not mention finality. Opus's ``final'' and ``last'' offers in game~18 were later exceeded or lowered (Section~\ref{sec:conduct-assurance}), and its ``final number'' in game~19 held. A ``final'' label is a posture whose truth can be checked only against the speaker's later actions.

\subsection{Assurances and changed plans: the \game{18} standing offer}

The engine has no contract object: a trade is a single offer that is accepted, rejected or countered, and nothing binds a player to terms stated in an earlier message. Assurances that span turns therefore exist only in language and memory, which makes them a clean test of what the records can establish about promises.

Late in game~18, Opus (seat~4) wants Sol's (seat~1) mortgaged Pacific Avenue, the deed that would complete its greens. At \decision{18}{419} Opus offers \$200; at \decision{18}{420} Sol counters at \$600, privately treating the deed as a blocker to sell only ``for a major liquidity swing''; at \decision{18}{421} Opus counters at \$250, publicly calling it ``my final offer''. Sol rejects at \decision{18}{422}. Over the next turns Opus repeatedly tells the table that the offer stands: ``Offer stands if you change your mind'' at \decision{18}{423}, ``My Pacific offer still stands'' at \decision{18}{425}, ``If you want cash for those reds, my Pacific offer still stands'' at \decision{18}{427}, and, after collecting a \$750 rent, ``Thanks for the \$750. Your roll, p1. My Pacific offer still stands'' at \decision{18}{431}. Opus's assistant content at that last decision adds that it had not raised the offer because ``extra cash now would only help p1 last longer''. Sol does not take the \$250 at first. It proposes a \$450 package at \decision{18}{437} and a \$350 package at \decision{18}{439}. In the interval Sol pays a second \$750 rent, sells its orange hotels and seven houses, and mortgages both remaining reds. At \decision{18}{441} it writes ``I'll take your standing offer: Pacific for \$250'', with a private note that refers explicitly to the earlier offer. At \decision{18}{442} Opus rejects: ``That offer's off the table now. Things have changed.'' Its returned reasoning cites Sol's \$49 of cash and the view that \$250 ``just extends their life briefly''. Later Opus offers \$150 and then \$225, again labelled ``final'' and ``still on the table'', and Sol answers at \$450, \$350 and \$300 without agreement.

The records support a specific account. The assurance was public, repeated and present in both parties' delivered memories at the invocation and the withdrawal. It was invoked explicitly and withdrawn after a material change in the invoker's finances, with that change stated as the reason. They do not support three stronger readings. Costly reliance is not established: Sol's intervening proposals asked for more than the standing price, and its sales and mortgages accompanied a second rent payment rather than any step toward the offer. Dishonesty at inception is not established: a feasible offer that is later withdrawn after the counterparty's position collapses is consistent with a sincere conditional offer. And the game's contract imposes no promise-keeping norm on bargaining. Work on commitments by language agents shows why these distinctions matter: whether a statement is a commitment, who may release it and when it expires must be defined before a withdrawal can be scored \citep{empty_commitments_v2,commitment_standing,shi2026cheaptalk}. Game~18 shows that the environment produces the raw material for such studies, namely cross-turn assurances with verifiable exposure, invocation and consequence. The chronology alone does not assign an obligation or release rule.

\subsection{Cooperation, opposition and third parties}

Four players competing for one win have reasons both to cooperate and to oppose the leader. Opposition appears in several reviewed games; explicit cooperation against a leader appears most clearly in game~4. The game~4 sequence above is the clearest conditional cooperation. One player publicly invites the others to form ``a multi-way or direct deal so p3 isn't the only one with a monopoly''. A second accepts the framing by name, pays the leader \$200 for the deed it needs, and the two negotiate a swap that gives each a buildable group. The cooperation was real in the sense that both parties obtained complete buildable groups; whether either gained against the leader was not tested. It was also bilateral and self-interested throughout: the false house-cost claim and the drop from a \$143 to a \$50 cash request were made inside it. Grok, the proposer, went on to win game~4.

Explicit refusals to fund a rival recur in several reviewed games. After Opus (seat~2) sells Illinois to Grok (seat~1) for \$500 in game~20, completing Grok's reds, Gemini (seat~4) ends its turn at \decision{20}{96} with a public appeal: ``let's not fund their houses by buying their spare properties for cash''. At \decision{20}{101} it rejects Grok's offer of Vermont for \$150, publicly and privately on the same ground. The appeal did not become an agreement. One of its addressees, Opus, later proposed cash purchases from Grok twice, neither of which settled. Gemini's rejections at \decision{9}{82} (``fund another house on your Dark Blues'') and \decision{16}{322} (``Keeping p2 cash-starved at \$142 is far better'') state the same reasoning: each rejects a proposed cash payment on the stated ground that it would finance a rival's next house or redemption. These are the agents' stated grounds. The rival-cash mechanisms in these two examples were not recomputed in the writing-stage check.

Denial can also misfire under the declared inheritance rule. At \decision{18}{335} Grok (seat~3) offers Gemini \$100 for mortgaged Illinois ``so p1 never completes the reds off you'', and Gemini accepts at \decision{18}{336} on the same stated ground. Eighteen turns later, at \decision{18}{360}, Grok owes Sol (seat~1) \$950 with \$64 in cash and declares bankruptcy. Its estate, including Illinois and both dark blues, passes to Sol, completing the reds the purchase was meant to block. Bankruptcy fixes the creditor by the debt. An estate transfer is therefore not a choice of beneficiary, and a denial purchase lasts only as long as its holder survives. A separate game~18 record shows a single received action completing the eventual winner's first group. At \decision{18}{83} Grok's received tool call accepts \$130 for Baltic, completing Opus's brown group, Opus's first monopoly in the game it went on to win. Grok's assistant content selected rejection, and its applied note says to ``Keep Baltic as a block''. This is a channel conflict with a material immediate consequence. Its contribution to the win was not tested, and the conflicting records do not support an intended-concession interpretation.

Across the adjudicated cases we found anti-leader appeals, denial purchases, refusals to fund rivals and conditional bilateral cooperation. We did not identify a multi-party coalition, cartel or deliberate kingmaking in these adjudicated cases. That is a bounded statement about the reviewed records, not proof of absence; the cohort has no private bilateral channel in which collusion could be arranged out of view, and the review did not code every trade response for leader-funding reasons. The eventual winner was the creditor in 54 of 60 eliminations (Section~\ref{sec:sol-economic-results}). Under a rule that sends the estate to whoever is owed, that concentration records which player held the debt, rather than an explicit choice of beneficiary. Studies that deliberately open collusive opportunities find behavior under designed incentives that this cohort did not provide \citep{shi2026emergentcollusion}.
